% Regression: t2_blocking — renders (a): A^{s}_{alpha,beta} = sum_gamma B^{s_1}_{alpha,gamma}.
% Formula: (a):  A^{s}_{alpha,beta} = sum_gamma B^{s_1}_{alpha,gamma}
% T2 blocking — RMP arXiv:2011.12127, fig. sec2fig11a1: fine-graining
% (inverse blocking) of a tensor-network site, 1D and 2D panels.
%
% Formula (a):  A^{s}_{alpha,beta} = sum_gamma B^{s_1}_{alpha,gamma}
%               C^{s_2}_{gamma,beta}  with s = (s_1, s_2):
% one MPS tensor splits into two joined by a new virtual bond.
%
% Formula (b):  the 2D analogue: one PEPS tensor A^s_{l,r,u,d} splits into
% a 2x2 patch of tensors contracted along the new internal bonds; the cut
% external bonds dangle as boundary legs and every site keeps its open
% physical index.
%
% Declarative reading:
%   (a) grid, physical=up: bead chains with open virtual stubs.
%   (b) tenkzlattice windows, frame=oblique: the RMP's sheared PEPS sheet;
%       boundary legs mark the virtual indices cut by the window and
%       physical=up gives the on-site vertical legs.  A 1x1 window IS the
%       single 5-leg tensor; a 2x2 window IS the fine-grained patch.
\documentclass[varwidth=19cm, border=6pt]{standalone}
\usepackage{amsmath, amssymb}
\usepackage{tenkz}
\begin{document}
(a)
\[
  % one 3-leg site  ->  two 3-leg sites sharing a new bond
  \begin{tenkz}[physical=up]
    \tn{}
  \end{tenkz}
  \;\longrightarrow\;
  \begin{tenkz}[physical=up]
    \tn{} & \tn{}
  \end{tenkz}
\]
(b)
\[
  % one 5-leg PEPS site (1x1 oblique window) -> 2x2 fine-grained patch
  \begin{tenkzlattice}[rows=1, cols=1, frame=oblique, boundary legs,
                       physical=up]
  \end{tenkzlattice}
  \;\longrightarrow\;
  \begin{tenkzlattice}[rows=2, cols=2, frame=oblique, boundary legs,
                       physical=up]
  \end{tenkzlattice}
\]
\end{document}
